%% methodology.tex
%%

%% ==============================
\section{Methodology}
\label{sec:Methodology}
%% ==============================

This section describes the experimental methodology used to reproduce the baseline paper's centralized Random Forest classifier and to build and evaluate the federated alternative. The methodology is designed to answer the research question: can a federated learning approach recover most of the centralized model's accuracy while addressing the privacy and scalability limitations the baseline paper identifies?

\subsection{Overview}
\label{subsec:method-overview}

Figure~\ref{fig:system-arch} illustrates the multi-site IoT QoS architecture. Six simulated IoT edge sites generate local QoS telemetry with non-IID traffic distributions. Figure~\ref{fig:training-comparison} contrasts the centralized training approach (pooling all sites' raw data) with the federated approach (training locally per site, combining predictions only).

% NOTE: TikZ diagrams commented out temporarily for successful compilation
% \begin{figure}[htbp]
% \centering
% \begin{tikzpicture}[
    node distance=1.5cm and 2cm,
    site/.style={rectangle, draw, fill=blue!20, minimum width=2.5cm, minimum height=1.2cm, align=center, rounded corners},
    data/.style={rectangle, draw, fill=green!20, minimum width=2cm, minimum height=0.8cm, align=center},
    model/.style={ellipse, draw, fill=orange!20, minimum width=2cm, minimum height=0.8cm, align=center},
    arrow/.style={->, >=stealth, thick}
]

% Title
\node[above] at (0,3.5) {\Large \textbf{Multi-Site IoT QoS Architecture}};

% Sites
\node[site] (site1) at (-6,2) {Site 1\\(Uniform)};
\node[site] (site2) at (-3,2) {Site 2\\(Burst)};
\node[site] (site3) at (0,2) {Site 3\\(Real-time)};
\node[site] (site4) at (3,2) {Site 4\\(Uniform)};
\node[site] (site5) at (6,2) {Site 5\\(Burst)};

% Local data
\node[data] (data1) at (-6,0.3) {Local\\Telemetry};
\node[data] (data2) at (-3,0.3) {Local\\Telemetry};
\node[data] (data3) at (0,0.3) {Local\\Telemetry};
\node[data] (data4) at (3,0.3) {Local\\Telemetry};
\node[data] (data5) at (6,0.3) {Local\\Telemetry};

% Local models
\node[model] (model1) at (-6,-1.5) {RF 1};
\node[model] (model2) at (-3,-1.5) {RF 2};
\node[model] (model3) at (0,-1.5) {RF 3};
\node[model] (model4) at (3,-1.5) {RF 4};
\node[model] (model5) at (6,-1.5) {RF 5};

% Federated ensemble
\node[model, fill=red!20, minimum width=4cm, align=center] (ensemble) at (0,-3.5) {Federated Ensemble\\(Probability Averaging)};

% Arrows
\draw[arrow] (site1) -- (data1);
\draw[arrow] (site2) -- (data2);
\draw[arrow] (site3) -- (data3);
\draw[arrow] (site4) -- (data4);
\draw[arrow] (site5) -- (data5);

\draw[arrow] (data1) -- (model1);
\draw[arrow] (data2) -- (model2);
\draw[arrow] (data3) -- (model3);
\draw[arrow] (data4) -- (model4);
\draw[arrow] (data5) -- (model5);

\draw[arrow] (model1) -- (ensemble);
\draw[arrow] (model2) -- (ensemble);
\draw[arrow] (model3) -- (ensemble);
\draw[arrow] (model4) -- (ensemble);
\draw[arrow] (model5) -- (ensemble);

% Note
\node[below, text width=10cm, align=center] at (0,-4.5) {\small \textit{Note: Each site trains locally on its own non-IID data.\\No raw telemetry leaves the site. Only model predictions are combined.}};

\end{tikzpicture}

% \caption{Multi-site IoT QoS architecture: six simulated edge sites, each with local telemetry and trained Random Forest model. The federated ensemble combines predictions via probability averaging without transmitting raw data.}
% \label{fig:system-arch}
% \end{figure}
% 
% \begin{figure}[htbp]
% \centering
% \begin{tikzpicture}[
    node distance=1.2cm and 1.5cm,
    data/.style={rectangle, draw, fill=green!20, minimum width=1.8cm, minimum height=0.6cm, text centered, font=\small},
    model/.style={ellipse, draw, fill=orange!20, minimum width=2cm, minimum height=0.8cm, text centered, font=\small},
    pool/.style={rectangle, draw, fill=yellow!20, minimum width=3cm, minimum height=1cm, text centered, font=\small},
    arrow/.style={->, >=stealth, thick}
]

% Centralized (Left)
\node[above] at (-4,4) {\Large \textbf{Centralized}};
\node[data] (c1) at (-5.5,2.5) {Site 1\\Data};
\node[data] (c2) at (-4.5,2.5) {Site 2\\Data};
\node[data] (c3) at (-3.5,2.5) {Site 3\\Data};
\node[data] (c4) at (-2.5,2.5) {Site 4\\Data};

\node[pool, align=center] (pool) at (-4,1) {Pooled Data\\+ SMOTE};
\node[model, fill=red!30, minimum width=2.5cm] (cmodel) at (-4,-0.5) {Centralized RF};

\draw[arrow] (c1) -- (pool);
\draw[arrow] (c2) -- (pool);
\draw[arrow] (c3) -- (pool);
\draw[arrow] (c4) -- (pool);
\draw[arrow] (pool) -- (cmodel);

% Federated (Right)
\node[above] at (4,4) {\Large \textbf{Federated}};
\node[data] (f1) at (2,2.5) {Site 1\\Data};
\node[data] (f2) at (3.5,2.5) {Site 2\\Data};
\node[data] (f3) at (5,2.5) {Site 3\\Data};
\node[data] (f4) at (6.5,2.5) {Site 4\\Data};

\node[pool, fill=green!30, align=center] (s1) at (2,1.2) {Local\\SMOTE};
\node[pool, fill=green!30, align=center] (s2) at (3.5,1.2) {Local\\SMOTE};
\node[pool, fill=green!30, align=center] (s3) at (5,1.2) {Local\\SMOTE};
\node[pool, fill=green!30, align=center] (s4) at (6.5,1.2) {Local\\SMOTE};

\node[model] (m1) at (2,-0.2) {RF 1};
\node[model] (m2) at (3.5,-0.2) {RF 2};
\node[model] (m3) at (5,-0.2) {RF 3};
\node[model] (m4) at (6.5,-0.2) {RF 4};

\node[model, fill=red!30, minimum width=3.5cm, align=center] (ens) at (4.25,-1.5) {Ensemble\\(Avg Probabilities)};

\draw[arrow] (f1) -- (s1);
\draw[arrow] (f2) -- (s2);
\draw[arrow] (f3) -- (s3);
\draw[arrow] (f4) -- (s4);

\draw[arrow] (s1) -- (m1);
\draw[arrow] (s2) -- (m2);
\draw[arrow] (s3) -- (m3);
\draw[arrow] (s4) -- (m4);

\draw[arrow] (m1) -- (ens);
\draw[arrow] (m2) -- (ens);
\draw[arrow] (m3) -- (ens);
\draw[arrow] (m4) -- (ens);

% Privacy note
\node[below, text width=4cm, align=center, font=\small] at (-4,-1.5) {\textcolor{red}{\textbf{Privacy Risk:}\\Raw data pooled}};
\node[below, text width=4cm, align=center, font=\small] at (4.25,-2.5) {\textcolor{green!60!black}{\textbf{Privacy Preserved:}\\Data never leaves sites}};

\end{tikzpicture}

% \caption{Centralized vs. federated training: centralized pools raw data (privacy risk), federated trains locally and combines predictions (privacy preserved).}
% \label{fig:training-comparison}
% \end{figure}

\subsection{Synthetic Multi-Site QoS Telemetry}
\label{subsec:telemetry}

Real multi-site OneM2M telemetry was not available, so a synthetic generator (\texttt{src/data/qos\_simulator.py}) produces QoS records for 6 simulated IoT edge sites. Each record contains six features drawn from the baseline paper's own QoS telemetry schema:
\begin{itemize}
  \item \textbf{RTT (Round-Trip Time)}: milliseconds, primary traffic-oriented metric
  \item \textbf{CPU utilization}: percentage (0-100\%)
  \item \textbf{RAM utilization}: percentage (0-100\%)
  \item \textbf{Success rate}: percentage (0-100\%), proportion of successful requests
  \item \textbf{Latency}: milliseconds, application-layer delay
  \item \textbf{Throughput}: requests/second, system capacity
\end{itemize}

Each site's traffic is a mixture of the baseline paper's own three scenarios -- \textit{uniform} (steady-state traffic), \textit{burst} (periodic load spikes), and \textit{real-time} (time-sensitive requests with tight latency bounds) -- but with a different, skewed mixture per site: 80\% dominated by one scenario, 10\% each of the other two. This skew is what makes this a genuine non-IID federated learning setting: no single site sees a representative sample of the overall traffic distribution. 1,500 records per site are generated per random seed, for a total of 9,000 records per seed across all six sites.

\subsubsection{Decision Labels and Thresholds}
Each record is labeled using the baseline paper's exact thresholds:
\begin{itemize}
  \item \textbf{Class 0 (Keep Local / Preferable)}: RTT $<$ 2s, success rate $>$ 95\%, CPU $<$ 70\%, RAM $<$ 70\%
  \item \textbf{Class 1 (Partial Offload / Acceptable)}: RTT between 2s and 4s, success rate 90-95\%, CPU 70-80\%, RAM 70-80\%
  \item \textbf{Class 2 (Full Offload / Critical)}: RTT $>$ 4s, success rate $<$ 90\%, CPU $>$ 80\%, or RAM $>$ 80\%
\end{itemize}

The synthetic generator applies these thresholds deterministically to produce class labels, then adds small amounts of Gaussian noise to features to simulate measurement error without changing labels. Class imbalance matches the baseline paper's reported distribution: approximately 48\% Class 0, 42\% Class 1, 10\% Class 2 (averaged across all sites).

\subsection{Training Pipeline}
\label{subsec:training}

For each of 5 random seeds (42, 43, 44, 45, 46), the following procedure is executed:

\begin{enumerate}
  \item \textbf{Data split}: Each site's 1,500 records are split 80/20 (stratified by label) into a local training set (1,200 records) and a local test set (300 records). The local test sets from all six sites are pooled into a single held-out test set (1,800 records total) used to evaluate all models.
  
  \item \textbf{Centralized training (baseline reproduction)}: All sites' training sets (6 $\times$ 1,200 = 7,200 records) are pooled into one centralized training set. SMOTE (Synthetic Minority Over-sampling Technique) is applied to balance the three classes to equal proportions. A single Random Forest classifier (200 trees, max depth 10, Gini impurity splitting criterion) is trained on the SMOTE-balanced pooled data.
  
  \item \textbf{Federated training (proposed improvement)}: For each of the 6 sites independently:
  \begin{itemize}
    \item Apply SMOTE locally to that site's training set (1,200 records) to balance the three classes within that site.
    \item Train a Random Forest classifier (200 trees, max depth 10, Gini splitting) on the locally SMOTE-balanced data.
    \item Store the trained model; no data or model parameters are shared with any other site or central server at this stage.
  \end{itemize}
  At inference time, each of the 6 site models predicts class probabilities for a given test instance; these 6 probability vectors are averaged element-wise, and the class with the highest averaged probability is returned as the ensemble prediction.
  
  \item \textbf{Single-site local-only training (lower bound)}: For each of the 6 sites independently, the same locally SMOTE-balanced Random Forest (200 trees, max depth 10) is trained as in the federated case, but evaluated in isolation on the pooled test set with no combination with any other site's model. This represents the performance a site achieves if it can neither centralize data nor federate with other sites.
\end{enumerate}

All three variants (centralized, federated, single-site local-only) use identical Random Forest hyperparameters. The only difference is the data-sharing strategy: centralized pools raw data; federated trains locally and combines predictions; single-site uses only its own local model. This design isolates the effect of the data-sharing strategy on model performance.

\subsection{Evaluation Metrics}
\label{subsec:metrics}

All models are evaluated on the same pooled, held-out test set (20\% of every site's data, stratified by label, never used in training by any model). Three metrics are reported:

\begin{itemize}
  \item \textbf{Accuracy}: overall proportion of correct predictions across all three classes.
  \item \textbf{Macro-F1 score}: unweighted average of per-class F1 scores, giving equal importance to all three classes regardless of their frequency in the test set. This metric is preferred over weighted-F1 or micro-F1 because it does not let the majority classes dominate the evaluation.
  \item \textbf{Critical Recall (class 2 recall)}: recall specifically on the ``Full Offload'' (Critical) class. This is the minority class the baseline paper needed SMOTE to handle (9.64\% of their real data), and the one most consequential to miss in production: a false negative on Critical means an overloaded system incorrectly stays local, risking degraded user experience or system failure.
\end{itemize}

Results are reported as mean $\pm$ standard deviation across 5 seeds. Statistical significance of differences between variants is not formally tested due to the small number of seeds, but consistency across seeds is noted qualitatively.

\subsection{Implementation Details}
\label{subsec:implementation}

All models are implemented in Python 3 using scikit-learn 1.3 for Random Forest and imbalanced-learn 0.11 for SMOTE. Training and evaluation are performed on a single machine (Intel Xeon CPU, 32GB RAM, no GPU required). Total runtime for all three variants across 5 seeds is approximately 15 minutes. Trained models are serialized via joblib for deployment; the federated ensemble's 6 per-site models total 42MB on disk (7MB per model), comparable to the centralized model's single 8MB file.

All code, synthetic telemetry generator, training scripts, and results are available in \texttt{iot-reliability/}.
